% ROZDZIAŁ 3 — plik włączany przez main.tex
\chapter{Miary zewnętrzne i konstrukcja Carathéodory'ego}
\label{ch:caratheodory}

\section{Problem konstrukcji miary}

W rozdziale \ref{ch:miary-wlasnosci} badaliśmy miarę już określoną na
$\sigma$-algebrze. Teraz odwrócimy pytanie: jeśli znamy ``rozmiary''
prostych zbiorów, to jak zbudować miarę na znacznie większej rodzinie?
Przykładowo na prostej chcemy przypisać przedziałowi $[a,b)$ długość
$b-a$, a następnie mierzyć zbiory otrzymane przez przeliczalne operacje.
Sam warunek skończonej addytywności jest za słaby. Kluczową rolę odgrywa
przeliczalna addytywność oraz możliwość pokrywania dowolnego zbioru
zbiorami prostszymi.

W całym rozdziale $X$ oznacza ustalony zbiór, a wszystkie funkcje
zbiorów przyjmują wartości w $[0,\infty]$, jeśli wyraźnie nie
zaznaczono inaczej. Przyjmujemy $\inf\varnothing=\infty$ oraz
$\sum_{n=1}^\infty a_n=\sup_{N\ge1}\sum_{n=1}^N a_n$ dla $a_n\ge0$.

\section{Miara zewnętrzna: definicja i własności}

\begin{definicja}[Miara zewnętrzna]
\label{def:outer-measure}
Funkcję $\mu^*:\mathcal P(X)\to[0,\infty]$ nazywamy
\emph{miarą zewnętrzną}, jeśli:
\begin{enumerate}[label=\textup{(\roman*)}]
\item $\mu^*(\varnothing)=0$;
\item $A\subseteq B$ implikuje $\mu^*(A)\le\mu^*(B)$;
\item dla dowolnego ciągu $(A_n)_{n\ge1}$ podzbiorów $X$ zachodzi
\begin{equation}\label{eq:outer-subadd}
\mu^*\Big(\bigcup_{n=1}^\infty A_n\Big)
\le\sum_{n=1}^\infty\mu^*(A_n).
\end{equation}
\end{enumerate}
Nie zakładamy, że $\mu^*$ jest miarą na $\mathcal P(X)$:
na ogół nie jest addytywna nawet dla rozłącznych zbiorów.
\end{definicja}

\begin{przyklad}[Najprostsze miary zewnętrzne]
Dla każdego zbioru $X$ funkcja $\mu^*(A)=0$ jest miarą zewnętrzną.
Miara licząca $\#(A)$ z przykładu \ref{ex:counting}, traktowana
jako funkcja na $\mathcal P(X)$, także jest miarą zewnętrzną.
Bardziej pouczający przykład otrzymujemy na $X=\{a,b\}$, gdy
$\mu^*(\varnothing)=0$ i $\mu^*(A)=1$ dla każdego $A\ne\varnothing$.
Wtedy $\mu^*(\{a\}\cup\{b\})=1\ne2=\mu^*(\{a\})+\mu^*(\{b\})$.
Wkrótce zobaczymy, że jej zbiory mierzalne w sensie Carathéodory'ego
to jedynie $\varnothing$ i $X$.
\end{przyklad}

\subsection{Miara zewnętrzna z kosztu pokryć}

\begin{twierdzenie}[Konstrukcja przez pokrycia]
\label{thm:outer-from-cover}
Niech $\mathcal C\subseteq\mathcal P(X)$ będzie rodziną zawierającą
$\varnothing$ oraz taką, że $X$ jest sumą przeliczalnie wielu zbiorów
z $\mathcal C$. Niech $c:\mathcal C\to[0,\infty]$ spełnia
$c(\varnothing)=0$. Dla każdego $E\subseteq X$ określmy
\begin{equation}\label{eq:outer-cover}
\mu_c^*(E):=\inf\Big\{\sum_{n=1}^\infty c(C_n):
C_n\in\mathcal C,\quad E\subseteq\bigcup_{n=1}^\infty C_n\Big\}.
\end{equation}
Wówczas $\mu_c^*$ jest miarą zewnętrzną.
\end{twierdzenie}

\begin{proof}
Pusty zbiór można pokryć ciągiem zbiorów pustych, więc
$0\le\mu_c^*(\varnothing)\le0$. Jeżeli $E\subseteq F$, to każde
pokrycie $F$ dopuszczalne w \eqref{eq:outer-cover} jest pokryciem $E$;
przy infimum po większej rodzinie nie można uzyskać większej wartości.

Pozostaje podaddytywność. Niech $E=\bigcup_{j\ge1}E_j$.
Jeżeli $\sum_j\mu_c^*(E_j)=\infty$, teza jest automatyczna.
W przeciwnym przypadku, przy ustalonym $\varepsilon>0$, dla każdego
$j$ wybieramy pokrycie $E_j\subseteq\bigcup_{k\ge1}C_{j,k}$, dla
którego
$\sum_{k\ge1}c(C_{j,k})<\mu_c^*(E_j)+\varepsilon2^{-j}$.
Taki wybór jest możliwy z definicji infimum; sumy skończone też
można uzupełnić pustymi zbiorami. Rodzina
$\{C_{j,k}:j,k\ge1\}$ jest przeliczalna i pokrywa $E$. Zatem
\[
\mu_c^*(E)\le \sum_{j,k\ge1}c(C_{j,k})
\le\sum_{j\ge1}\mu_c^*(E_j)+\varepsilon.
\]
Ponieważ $\varepsilon>0$ było dowolne, otrzymujemy
\eqref{eq:outer-subadd}. Przestawianie sum jest tu dozwolone,
gdyż wszystkie składniki są nieujemne: obie sumy to kresy górne
sum po skończonych podzbiorach $\mathbb N^2$.
\end{proof}

\begin{uwaga}[Rola założeń]
W konstrukcji wystarczy, aby każde $E\subseteq X$ miało jakieś
przeliczalne pokrycie zbiorami z $\mathcal C$; przyjęte założenie
zapewnia to jednocześnie wszystkim $E$. Nie wymagamy addytywności $c$.
Z drugiej strony samo pokrywanie nie gwarantuje, że
$\mu_c^*(C)=c(C)$ dla $C\in\mathcal C$.
Równość uzyskamy dopiero przy założeniu, że $c$ jest premiarą.
\end{uwaga}

\section{Kryterium mierzalności Carathéodory'ego}

\begin{definicja}[Zbiór mierzalny w sensie Carathéodory'ego]
\label{def:caratheodory-measurable}
Niech $\mu^*$ będzie miarą zewnętrzną na $X$. Zbiór $E\subseteq X$
nazywamy $\mu^*$-\emph{mierzalnym}, jeśli dla każdego $T\subseteq X$
zachodzi
\begin{equation}\label{eq:caratheodory-criterion}
\mu^*(T)=\mu^*(T\cap E)+\mu^*(T\setminus E).
\end{equation}
Rodzinę takich zbiorów oznaczamy przez $\mathcal M_{\mu^*}$.
\end{definicja}

Z nierówności \eqref{eq:outer-subadd} zawsze wynika
$\mu^*(T)\le\mu^*(T\cap E)+\mu^*(T\setminus E)$.
Weryfikacja \eqref{eq:caratheodory-criterion} polega więc na
udowodnieniu \emph{odwrotnej} nierówności, także wtedy, gdy
$\mu^*(T)=\infty$.

\begin{twierdzenie}[Twierdzenie Carathéodory'ego]
\label{thm:caratheodory}
Jeśli $\mu^*$ jest miarą zewnętrzną na $X$, to
$\mathcal M_{\mu^*}$ jest $\sigma$-algebrą na $X$, zaś
$\mu:=\mu^*|_{\mathcal M_{\mu^*}}$ jest miarą. Ponadto miara
$\mu$ jest zupełna: każdy podzbiór zbioru $\mu$-miary zero jest
$\mu^*$-mierzalny i ma miarę zero.
\end{twierdzenie}

\begin{proof}
\emph{Krok 1: dopełnienia i zbiory trywialne.}
Dla $E=\varnothing$ równość
\eqref{eq:caratheodory-criterion} ma postać
$\mu^*(T)=0+\mu^*(T)$. Warunek jest symetryczny względem
$E$ i $X\setminus E$, toteż mierzalność $E$ pociąga mierzalność
jego dopełnienia. W szczególności $X\in\mathcal M_{\mu^*}$.

\emph{Krok 2: sumy dwóch zbiorów.}
Dla $E,F\in\mathcal M_{\mu^*}$ i dowolnego $T\subseteq X$
rozkładamy najpierw $T$ według $E$, a potem $T\setminus E$
według $F$:
\begin{align*}
\mu^*(T)
&=\mu^*(T\cap E)+\mu^*(T\setminus E)\\
&=\mu^*(T\cap E)+\mu^*((T\setminus E)\cap F)
  +\mu^*(T\setminus(E\cup F))\\
&\ge\mu^*(T\cap(E\cup F))
  +\mu^*(T\setminus(E\cup F)).
\end{align*}
Ostatnia nierówność wynika z podaddytywności, gdyż
$T\cap(E\cup F)=(T\cap E)\cup((T\setminus E)\cap F)$.
Nierówność w drugą stronę obowiązuje dla każdego zbioru, a zatem
$E\cup F$ jest mierzalny. Mamy też zamkniętość na skończone
przekroje, dzięki prawom de Morgana.

\emph{Krok 3: przeliczalne sumy rozłączne.}
Niech $(E_n)_{n\ge1}$ będzie parami rozłącznym ciągiem
zbiorów mierzalnych, a $E=\bigcup_nE_n$.
Wielokrotnie stosując kryterium do $E_1,\ldots,E_N$,
otrzymujemy dla dowolnego $T\subseteq X$
\[
\mu^*(T)=\sum_{n=1}^N\mu^*(T\cap E_n)
+\mu^*\Big(T\setminus\bigcup_{n=1}^N E_n\Big).
\]
Ponieważ $T\setminus E\subseteq T\setminus\bigcup_{n=1}^N E_n$,
monotoniczność daje
$\mu^*(T)\ge\sum_{n=1}^N\mu^*(T\cap E_n)
+\mu^*(T\setminus E)$.
Przechodząc do kresu górnego po $N$, otrzymujemy
\[
\mu^*(T)\ge\sum_{n=1}^\infty\mu^*(T\cap E_n)
+\mu^*(T\setminus E)
\ge\mu^*(T\cap E)+\mu^*(T\setminus E).
\]
Ostatnie przejście korzysta z podaddytywności dla
$T\cap E=\bigcup_n(T\cap E_n)$. To dowodzi mierzalności $E$.

\emph{Krok 4: dowolne sumy przeliczalne.}
Dla mierzalnych $A_n$ stosujemy rozłącznienie
$E_1=A_1$ oraz $E_n=A_n\setminus\bigcup_{k<n}A_k$ dla $n\ge2$.
Na mocy poprzednich kroków wszystkie $E_n$ są mierzalne i
$\bigcup_n A_n=\bigcup_n E_n$; zatem
$\mathcal M_{\mu^*}$ jest $\sigma$-algebrą.

\emph{Krok 5: przeliczalna addytywność.}
Dla rozłącznych mierzalnych $E_n$ podstawiamy $T=E=\bigcup_nE_n$
do nierówności z kroku 3. Otrzymujemy
$\mu^*(E)\ge\sum_n\mu^*(E_n)$, a odwrotną nierówność zapewnia
\eqref{eq:outer-subadd}. Ponadto $\mu^*(\varnothing)=0$,
więc ograniczenie $\mu$ jest miarą.

\emph{Krok 6: zupełność.}
Niech $N\in\mathcal M_{\mu^*}$, $\mu^*(N)=0$ i $S\subseteq N$.
Dla dowolnego $T$ z monotoniczności mamy
$\mu^*(T\cap S)=0$ i
$\mu^*(T\setminus S)\le\mu^*(T)$.
Stąd
$\mu^*(T\cap S)+\mu^*(T\setminus S)\le\mu^*(T)$.
Podaddytywność daje nierówność przeciwną. Zatem $S$ jest
mierzalny oraz $\mu^*(S)=0$.
\end{proof}

\begin{przyklad}[Miara zewnętrzna, która nie mierzy wszystkich zbiorów]
Niech $X=\{a,b\}$ oraz $\mu^*(A)=1$ dla $A\ne\varnothing$.
Gdyby $E=\{a\}$ było mierzalne, podstawienie $T=X$ dałoby
$1=\mu^*(X)=\mu^*(\{a\})+\mu^*(\{b\})=2$,
sprzeczność. Stąd $\mathcal M_{\mu^*}=\{\varnothing,X\}$.
To pokazuje, że kryterium Carathéodory'ego ma rzeczywistą treść.
\end{przyklad}

\section{Premiary na algebrach zbiorów}

\begin{definicja}[Premiara]
\label{def:premeasure}
Niech $\mathcal A$ będzie algebrą podzbiorów $X$.
Funkcję $\mu_0:\mathcal A\to[0,\infty]$ nazywamy
\emph{premiarą}, jeżeli $\mu_0(\varnothing)=0$ i dla każdego
ciągu parami rozłącznych $A_n\in\mathcal A$ takiego, że
$A=\bigcup_{n\ge1}A_n\in\mathcal A$, zachodzi
\begin{equation}\label{eq:premeasure}
\mu_0(A)=\sum_{n=1}^\infty\mu_0(A_n).
\end{equation}
Warunek nie zakłada, że $\bigcup_nA_n$ zawsze należy do
$\mathcal A$; algebra nie musi być $\sigma$-algebrą.
\end{definicja}

\begin{lemat}[Podaddytywność premiary na pokryciach]
\label{lem:premeasure-cover}
Jeżeli $\mu_0$ jest premiarą na $\mathcal A$, zaś
$A,A_n\in\mathcal A$ i $A\subseteq\bigcup_{n\ge1}A_n$,
to $\mu_0(A)\le\sum_{n\ge1}\mu_0(A_n)$.
\end{lemat}

\begin{proof}
Tworzymy rozłączne zbiory
$B_1=A\cap A_1$ i
$B_n=(A\cap A_n)\setminus\bigcup_{k<n}A_k$ dla $n\ge2$.
Wszystkie należą do algebry, a $A=\bigsqcup_nB_n$.
Z \eqref{eq:premeasure} oraz monotoniczności premiary
(wynikającej ze skończonej addytywności) dostajemy
$\mu_0(A)=\sum_n\mu_0(B_n)\le\sum_n\mu_0(A_n)$.
\end{proof}

\begin{przyklad}[Skończona addytywność nie wystarcza]
Na $X=\mathbb N$ niech $\mathcal A$ będzie algebrą
zbiorów skończonych i kofinitywnych. Określmy
$\nu_0(A)=0$, gdy $A$ jest skończony, oraz
$\nu_0(A)=1$, gdy $\mathbb N\setminus A$ jest skończony.
Jest to funkcja skończenie addytywna: dwa rozłączne
zbiory kofinitywne nie istnieją, a dołączenie zbioru
skończonego nie zmienia wartości. Nie jest premiarą, gdyż
$\mathbb N=\bigsqcup_{n\ge1}\{n\}$, natomiast
$\nu_0(\mathbb N)=1\ne\sum_n\nu_0(\{n\})=0$.
Każda miara rozszerzająca $\nu_0$ musiałaby spełniać
sprzeczną równość, toteż takie rozszerzenie nie istnieje.
\end{przyklad}

\begin{twierdzenie}[Kryterium ciągłości premiary skończonej]
\label{thm:premeasure-continuity}
Niech $\mathcal A$ będzie algebrą podzbiorów $X$, a
$\mu_0:\mathcal A\to[0,\infty)$ funkcją skończenie addytywną
spełniającą $\mu_0(X)<\infty$ i $\mu_0(\varnothing)=0$.
Wówczas $\mu_0$ jest premiarą wtedy i tylko wtedy, gdy
$A_n\in\mathcal A$, $A_n\downarrow\varnothing$ implikuje
$\mu_0(A_n)\to0$.
\end{twierdzenie}

\begin{proof}
Jeśli $\mu_0$ jest premiarą, to dla
$A_n\downarrow\varnothing$ zbiory $D_n=A_n\setminus A_{n+1}$
są rozłączne i $A_1=\bigsqcup_nD_n$.
W konsekwencji
$\mu_0(A_n)=\sum_{k\ge n}\mu_0(D_k)\to0$,
gdyż szereg ma skończoną sumę $\mu_0(A_1)\le\mu_0(X)$.

Odwrotnie, niech $A=\bigsqcup_n A_n$ i wszystkie te zbiory
należą do $\mathcal A$. Dla
$R_N=A\setminus\bigcup_{n=1}^N A_n$ mamy
$R_N\downarrow\varnothing$, a skończona addytywność daje
$\mu_0(A)=\sum_{n=1}^N\mu_0(A_n)+\mu_0(R_N)$.
Po przejściu do granicy otrzymujemy \eqref{eq:premeasure}.
\end{proof}

\begin{uwaga}
Nie wolno bez zmian stosować twierdzenia
\ref{thm:premeasure-continuity} przy $\mu_0(X)=\infty$.
Na przykład miara licząca na $\mathbb N$ jest premiarą,
lecz dla $A_n=\{n,n+1,\ldots\}$ mamy $A_n\downarrow\varnothing$
i $\mu_0(A_n)=\infty$ dla wszystkich $n$.
\end{uwaga}

\section{Twierdzenie o rozszerzaniu miar}

\begin{twierdzenie}[Rozszerzenie Carathéodory'ego]
\label{thm:extension-caratheodory}
Niech $\mathcal A$ będzie algebrą podzbiorów $X$, a
$\mu_0$ premiarą na $\mathcal A$. Zdefiniujmy $\mu^*$
przez \eqref{eq:outer-cover}, biorąc $\mathcal C=\mathcal A$
i $c=\mu_0$. Wówczas:
\begin{enumerate}[label=\textup{(\roman*)}]
\item $\mu^*$ jest miarą zewnętrzną;
\item każdy $A\in\mathcal A$ jest $\mu^*$-mierzalny;
\item $\mu^*(A)=\mu_0(A)$ dla wszystkich $A\in\mathcal A$;
\item $\mu=\mu^*|_{\mathcal M_{\mu^*}}$ jest zupełną miarą
rozszerzającą $\mu_0$, a jej ograniczenie do
$\sigma(\mathcal A)$ jest miarą na $\sigma(\mathcal A)$.
Miara na $\sigma(\mathcal A)$ nie musi jednak być zupełna.
\end{enumerate}
\end{twierdzenie}

\begin{proof}
Punkt (i) wynika z twierdzenia \ref{thm:outer-from-cover},
gdyż $X\in\mathcal A$.
Ustalmy $A\in\mathcal A$ i dowolny $T\subseteq X$.
Dla każdego pokrycia $T\subseteq\bigcup_nC_n$ przez
$C_n\in\mathcal A$ zbiory $C_n\cap A$ pokrywają $T\cap A$,
a $C_n\setminus A$ pokrywają $T\setminus A$. Dlatego
\begin{align*}
\mu^*(T\cap A)+\mu^*(T\setminus A)
&\le\sum_n\mu_0(C_n\cap A)+\sum_n\mu_0(C_n\setminus A)\\
&=\sum_n\mu_0(C_n).
\end{align*}
Równość korzysta ze skończonej addytywności $\mu_0$.
Biorąc infimum po pokryciach, otrzymujemy
$\mu^*(T\cap A)+\mu^*(T\setminus A)\le\mu^*(T)$.
W połączeniu z podaddytywnością jest to warunek
\eqref{eq:caratheodory-criterion}; dowodzi (ii).

Dla każdego $A\in\mathcal A$ pokrycie $A$ samym $A$
(i pozostałymi zbiorami pustymi) daje
$\mu^*(A)\le\mu_0(A)$. Z drugiej strony
lemat \ref{lem:premeasure-cover} zapewnia
$\mu_0(A)\le\sum_n\mu_0(C_n)$ dla dowolnego pokrycia
$A\subseteq\bigcup_nC_n$. Infimum daje
$\mu_0(A)\le\mu^*(A)$, a więc (iii).

Punkt (iv) jest już konsekwencją twierdzenia
\ref{thm:caratheodory}: $\mathcal M_{\mu^*}$ jest
$\sigma$-algebrą zawierającą $\mathcal A$, zatem
zawiera $\sigma(\mathcal A)$; ograniczenie $\mu^*$
jest na niej miarą, zaś na $\mathcal M_{\mu^*}$ miarą zupełną.
\end{proof}

\begin{definicja}[\texorpdfstring{$\sigma$}{sigma}-skończona premiara]\label{idx:03:definicja:09}

Premiarę $\mu_0$ na algebrze $\mathcal A$ nazywamy
$\sigma$-\emph{skończoną}, jeżeli istnieją
$A_n\in\mathcal A$ takie, że $X=\bigcup_nA_n$ oraz
$\mu_0(A_n)<\infty$ dla każdego $n$.
Możemy zawsze zastąpić $A_n$ sumami skończonymi
$A_1\cup\cdots\cup A_n$, otrzymując ciąg rosnący.
\end{definicja}

\begin{twierdzenie}[Jednoznaczność rozszerzenia]
\label{thm:extension-uniqueness}
Jeśli premiara $\mu_0$ na algebrze $\mathcal A$ jest
$\sigma$-skończona, to istnieje dokładnie jedna miara
na $\sigma(\mathcal A)$, której ograniczenie do
$\mathcal A$ jest równe $\mu_0$.
\end{twierdzenie}

\begin{proof}
Istnienie daje twierdzenie \ref{thm:extension-caratheodory}.
Niech $\mu$ i $\nu$ będą dwoma rozszerzeniami.
Wybierzmy rosnący ciąg $A_n\in\mathcal A$ o
$\bigcup_n A_n=X$ i $\mu_0(A_n)<\infty$.
Dla ustalonego $n$ funkcje
$\mu_n(E)=\mu(E\cap A_n)$ oraz $\nu_n(E)=\nu(E\cap A_n)$,
$E\in\sigma(\mathcal A)$, są miarami skończonymi.
Zgadzają się na algebrze $\mathcal A$, ponieważ
$E\cap A_n\in\mathcal A$ dla $E\in\mathcal A$.
Algebra jest układem $\pi$ i generuje
$\sigma(\mathcal A)$. Z twierdzenia
\ref{thm:uniqueness-finite} otrzymujemy
$\mu_n(E)=\nu_n(E)$ dla każdego $E\in\sigma(\mathcal A)$.
Ponieważ $E\cap A_n\uparrow E$, ciągłość miary z dołu
(twierdzenie \ref{thm:measure-continuity}) daje
$\mu(E)=\lim_n\mu_n(E)=\lim_n\nu_n(E)=\nu(E)$.
\end{proof}

\begin{przyklad}[Dlaczego potrzebna jest $\sigma$-skończoność]
\label{ex:nonunique-extension}
Na $X=\mathbb R$ niech $\mathcal A$ będzie algebrą
generowaną przez przedziały $[a,b)$, gdzie $a<b$ są
rzeczywiste. Każdy niepusty element $\mathcal A$ zawiera
niezdegenerowany przedział; ponadto
$\sigma(\mathcal A)=\mathcal B(\mathbb R)$.
Na zbiorach borelowskich rozważmy dwie miary:
\[
\mu(E)=\#E,\qquad
\nu(E)=\begin{cases}0,&E\text{ jest przeliczalny},\\
\infty,&E\text{ jest nieprzeliczalny}.
\end{cases}
\]
Funkcja $\nu$ jest miarą, gdyż przeliczalna suma zbiorów
przeliczalnych jest przeliczalna, a suma zawierająca
zbiór nieprzeliczalny jest nieprzeliczalna. Na każdym
niepustym elemencie $\mathcal A$ obie miary mają wartość
$\infty$, zaś na pustym obie są równe zero. Są to więc
różne rozszerzenia tej samej premiary:
$\mu(\{0\})=1$, natomiast $\nu(\{0\})=0$.
Premiara nie jest $\sigma$-skończona: jej jedynym
zbiorem o skończonej mierze jest $\varnothing$.
\end{przyklad}

\begin{uwaga}[Różne role założeń]
Premiara daje istnienie rozszerzenia na $\sigma(\mathcal A)$
bez założenia $\sigma$-skończoności. To założenie pojawia się
dopiero przy jednoznaczności rozszerzenia oraz przy
identyfikacji całej rodziny zbiorów Carathéodory'ego
z uzupełnieniem $\sigma(\mathcal A)$. W szczególności
stwierdzenie o istnieniu i stwierdzenie o jednoznaczności
nie są jednym i tym samym twierdzeniem.
\end{uwaga}

\section{Kiedy mierzalność Carathéodory'ego jest uzupełnieniem rozszerzenia?}

Rozszerzenie otrzymane z twierdzenia
\ref{thm:extension-caratheodory} jest określone na
$\mathcal M_{\mu^*}$, zwykle większej niż $\sigma(\mathcal A)$.
Nie należy utożsamiać tych rodzin bez dowodu.
Poniższy rezultat określa ich związek przy
$\sigma$-skończoności.

\begin{lemat}[Aproksymacja od góry]
\label{lem:outer-approx-algebra}
Niech $\mu_0$ będzie premiarą na algebrze $\mathcal A$ i
$\mu^*$ jej miarą zewnętrzną. Jeśli $E\subseteq X$ oraz
$\mu^*(E)<\infty$, to dla każdego $\varepsilon>0$
istnieje $B\in\sigma(\mathcal A)$, $E\subseteq B$, takie, że
$\mu^*(B)\le\mu^*(E)+\varepsilon$.
Jeżeli dodatkowo $E\in\mathcal M_{\mu^*}$, to
$\mu(B\setminus E)<\varepsilon$.
\end{lemat}

\begin{proof}
Wybieramy pokrycie $E\subseteq\bigcup_nA_n$ przez
$A_n\in\mathcal A$ z $\sum_n\mu_0(A_n)<\mu^*(E)+\varepsilon$.
Wystarczy wziąć $B=\bigcup_nA_n$. Z podaddytywności i
równości $\mu^*(A_n)=\mu_0(A_n)$ dostajemy
$\mu^*(B)\le\sum_n\mu_0(A_n)<\mu^*(E)+\varepsilon$.
Jeśli $E$ jest mierzalny, to $B\setminus E$ jest mierzalny
i $\mu(B)=\mu(E)+\mu(B\setminus E)$.
Obie pierwsze wartości są skończone, więc
$\mu(B\setminus E)<\varepsilon$.
\end{proof}

\begin{twierdzenie}[Uzupełnienie rozszerzenia]
\label{thm:caratheodory-completion}
Załóżmy, że premiara $\mu_0$ na algebrze $\mathcal A$
jest $\sigma$-skończona. Niech $\mu^*$ będzie wyznaczoną
przez nią miarą zewnętrzną, a $\mu_B=\mu^*|_{\sigma(\mathcal A)}$.
Wówczas przestrzeń
$(X,\mathcal M_{\mu^*},\mu^*|_{\mathcal M_{\mu^*}})$
jest uzupełnieniem przestrzeni
$(X,\sigma(\mathcal A),\mu_B)$ w sensie twierdzenia
\ref{thm:completion}. Dokładniej: dla każdego
$E\in\mathcal M_{\mu^*}$ istnieją
$C,D\in\sigma(\mathcal A)$ takie, że
$C\subseteq E\subseteq D$ i $\mu_B(D\setminus C)=0$.
\end{twierdzenie}

\begin{proof}
Wybierzmy rosnący ciąg $A_n\in\mathcal A$ pokrywający $X$,
przy czym $\mu_0(A_n)<\infty$. Dla ustalonego mierzalnego
$E$ zbiór $E_n=E\cap A_n$ jest mierzalny i ma skończoną miarę.
Dla każdego $k,n\ge1$ lemat \ref{lem:outer-approx-algebra}
daje $B_{n,k}\in\sigma(\mathcal A)$ z
$E_n\subseteq B_{n,k}$ oraz
$\mu(B_{n,k}\setminus E_n)<2^{-n}/k$.
Połóżmy $D_k=\bigcup_{n\ge1}B_{n,k}$.
Wtedy $D_k\in\sigma(\mathcal A)$,
$E\subseteq D_k$, gdyż $E=\bigcup_nE_n$, i
\[
\mu(D_k\setminus E)
\le\sum_n\mu(B_{n,k}\setminus E)
\le\sum_n\mu(B_{n,k}\setminus E_n)<\frac1k.
\]
Zbiór $D=\bigcap_{k\ge1}D_k$ należy do
$\sigma(\mathcal A)$, zawiera $E$ i spełnia
$\mu(D\setminus E)=0$, ponieważ
$D\setminus E\subseteq D_k\setminus E$ dla każdego $k$.
Powtarzamy konstrukcję dla mierzalnego zbioru
$X\setminus E$, otrzymując zbiór należący do
$\sigma(\mathcal A)$, mianowicie $H\supseteq X\setminus E$,
taki że $\mu(H\setminus(X\setminus E))=0$.
Wtedy $C=X\setminus H\in\sigma(\mathcal A)$,
$C\subseteq E$ i $\mu(E\setminus C)=0$.
Wobec $C\subseteq E\subseteq D$ mamy
$D\setminus C=(D\setminus E)\sqcup(E\setminus C)$,
zatem $\mu_B(D\setminus C)=0$.

Każdy zbiór z uzupełnienia $\sigma(\mathcal A)$ różni
się od zbioru z $\sigma(\mathcal A)$ o podzbiór
zbioru $\mu_B$-zerowego, a więc jest
$\mu^*$-mierzalny z zupełności twierdzenia
\ref{thm:caratheodory}. Przed chwilą udowodniliśmy
zawieranie przeciwne, kończąc dowód.
\end{proof}

\begin{uwaga}[Dlaczego to rozróżnienie jest ważne?]
W rozdziale \ref{ch:lebesgue} rozróżnimy zbiory borelowskie
i zbiory mierzalne w sensie Lebesgue'a. Konstrukcja
Carathéodory'ego od razu daje miarę zupełną; ograniczenie
do borelowskiej $\sigma$-algebry nie musi być zupełne.
\end{uwaga}

\section{Jeszcze jeden sposób interpretacji konstrukcji}

Dla zbioru $E\subseteq X$ liczba $\mu^*(E)$ jest
najmniejszym kosztem \emph{przeliczalnego pokrycia} $E$
przez zbiory z wyjściowej algebry. Mierzalność w sensie
Carathéodory'ego oznacza, że przecięcie z $E$ oraz z jego
dopełnieniem nie powoduje sztucznej utraty kosztu,
niezależnie od badanego zbioru $T$.

Następujące trzy kroki to logiczny szkielet rozdziału:
\begin{enumerate}
\item z premiary $\mu_0$ tworzymy $\mu^*$ na całej $\mathcal P(X)$;
\item warunek \eqref{eq:caratheodory-criterion} wyodrębnia
$\sigma$-algebrę, na której $\mu^*$ jest miarą;
\item przy $\sigma$-skończoności dostajemy jedyne
rozszerzenie na $\sigma(\mathcal A)$, a po uzupełnieniu
pełną dziedzinę mierzalności Carathéodory'ego.
\end{enumerate}
W następnym rozdziale zrealizujemy ten program dla
objętości prostopadłościanów w $\mathbb R^d$.